\par\bigskip\Needspace{9\baselineskip}
\pdfbookmark[2]{Second Method, [5]}{ceres-second-5}
\begin{center}[5.]\end{center}
\begingroup\small
\noindent\hspace*{\parindent}[Once all terms of expression 2) have been reduced in this way to the form
\[
A\sin(iD\pm\text{\AstroSymbols ⚳}\mp\AscendingNode+kM-k'M')
\]
the corresponding term is
\[
\begin{aligned}
&5)\qquad\text{term in }\frac{\theta}{\sin i}=\\
&\quad-A\cos V\int\sin V.\sin(iD\pm\text{\AstroSymbols ⚳}\mp\AscendingNode+kM-k'M')\,n\,\mathrm{d}t\\
&\quad+A\sin V\int\cos V.\sin(iD\pm\text{\AstroSymbols ⚳}\mp\AscendingNode+kM-k'M')\,n\,\mathrm{d}t.
\end{aligned}
\]

The next step is to expand \(\displaystyle\genfrac{}{}{0pt}{}{\sin}{\cos}V\) in multiples of \(M\). This, too, can be done by interpolation. The result was
\[
\begin{aligned}
\cos V&=0{,}07848+0{,}99308\cos M-0{,}07784\cos2M\quad\text{etc.}\\
\sin V&=\phantom{0{,}07848+}0{,}99461\sin M-0{,}07792\sin2M\quad\text{etc.}
\end{aligned}
\]

Once the products with these quantities under the integral signs in 5) have been expanded, the integration can be carried out directly. Multiplying again by \(\cos V\) and \(\sin V\), respectively, gives the corresponding term in \(\theta\). The term 3) used as an example above thus gave
\SourceLeaf{413}
\[
\begin{aligned}
&-27''{,}47\sin(3D+\text{\AstroSymbols ⚳}-\AscendingNode)\\
&-\phantom{2}4''{,}56\sin(3D+\text{\AstroSymbols ⚳}-\AscendingNode+M)\\
&+\phantom{2}2''{,}26\sin(3D+\text{\AstroSymbols ⚳}-\AscendingNode-M)\\
&+\phantom{2}2''{,}81\sin(3D+\text{\AstroSymbols ⚳}-\AscendingNode+M').
\end{aligned}
\]

The full numerical results are given in the following letter, which also shows how terms whose arguments differ only by constants were combined into a single term.]
\par\endgroup
\begin{center}\rule{.30\linewidth}{.4pt}\end{center}
